\documentclass[a4paper]{article}

\usepackage{amsmath}
\usepackage{amssymb}
\usepackage{array}
\usepackage{graphicx}
\usepackage{hyperref}
\usepackage{listings}
\usepackage{longtable}
\usepackage{natbib}
\usepackage{nicefrac}
\usepackage{ulem}

\newcommand{\pequod}{PEQUOD}
\newcommand{\pequodexpanded}{Parallel Quasi-Geostrophic Model}
\newcommand{\version}{0.9.5-dev}

\begin{document}

\title{Getting started with \pequod\ \version}
\maketitle

\section{Introduction}

This document gives step-by-step instructions for getting started with \pequod.
Specifically, these steps enable one to compile and run \pequod\ on a vanilla
installation of Ubuntu 12.04 (Precise Pangolin)\footnote{Available from
\url{http://www.ubuntu.com}.}. Detailed instructions concerning the
configuration and usage of \pequod\ can be found in the more comprehensive
manual.

\section{Dependencies}

\pequod\ has a small number of required dependencies. Packages required to
compile the model with the gfortran compiler can be installed via:
\begin{lstlisting}[language = bash, frame = single]
sudo apt-get install gfortran make libcloog-ppl0 \
  liblapack-dev libopenmpi-dev openmpi-bin
\end{lstlisting}
Packages required to build the manual can be installed via:
\begin{lstlisting}[language = bash, frame = single]
sudo apt-get install texlive texlive-latex-extra
\end{lstlisting}

\section{Building the manual}

In the \pequod\ root directory the manual is compiled with the command:
\begin{lstlisting}[language = bash, frame = single]
make manual
\end{lstlisting}
and can be viewed with the command:
\begin{lstlisting}[language = bash, frame = single]
evince manual/manual.pdf
\end{lstlisting}

\section{Running the model in serial}

The source files supplied with \pequod\ are configured to use the gfortran
compiler, with a low resolution three-layer wind-driven double gyre model
configuration. The model makefiles require that the MPI directory be specified
via the \verb+MPI_DIR+ environment variable. This can be set via the command:
\begin{lstlisting}[language = bash, frame = single]
export MPI_DIR=/usr
\end{lstlisting}
The model is compiled with (executed in the \pequod\ root directory):
\begin{lstlisting}[language = bash, frame = single]
make qg3
\end{lstlisting}
and can be run with:
\begin{lstlisting}[language = bash, frame = single]
./qg3
\end{lstlisting}

\section{Running the model in parallel}

The model can be run in parallel by specifying a domain decomposition in
\verb+qg3p_params.F90+ (see the manual for full details). As an example, if
one replaces the lines:
\begin{lstlisting}[language = bash, frame = single]
  ! Grid
  integer, parameter :: layers = 3
  integer, parameter, dimension(1) :: iis = (/257/)
  integer, parameter, dimension(1) :: jjs = (/257/)
  ! The maximum grid size over all processes
  integer, parameter :: ii_m = 257, jj_m = 257
\end{lstlisting}
with:
\begin{lstlisting}[language = bash, frame = single]
  ! Grid
  integer, parameter :: layers = 3
  integer, parameter, dimension(1) :: iis = (/257/)
  integer, parameter, dimension(2) :: jjs = (/128, 129/)
  ! The maximum grid size over all processes
  integer, parameter :: ii_m = 257, jj_m = 129
\end{lstlisting}
then after recompiling the model:
\begin{lstlisting}[language = bash, frame = single]
make qg3
\end{lstlisting}
the model can be run on two processes with:
\begin{lstlisting}[language = bash, frame = single]
$MPI_DIR/bin/mpirun -np 2 ./qg3
\end{lstlisting}

\section{Configuring the model}

The model is primarily configured via preprocessor definitions in the
\verb+config.h+ header file, and Fortran parameters in the
\verb+qg3p_params.F90+ source file. Full details of the configuration of the
model via these files are given in the manual. Note that the model must be
recompiled after any modification to these source files with:
\begin{lstlisting}[language = bash, frame = single]
make qg3
\end{lstlisting}

%\bibliography{bibliography}
%\bibliographystyle{plainnat}

\end{document}
